\documentclass[11pt]{article}
\usepackage[T1]{fontenc}
\usepackage{amsmath,amssymb,booktabs}
\usepackage[margin=1in]{geometry}
\usepackage[hidelinks]{hyperref}
\usepackage{tabularx}
\title{Proof Dossier: Spectral Residual-to-Scalar Loss Bridge}
\author{Project Koios Formal}
\date{}
\begin{document}
\maketitle
\section{Identity and status}
\begin{tabularx}{\textwidth}{@{}l>{\raggedright\arraybackslash}X@{}}
\toprule
Statement ID & \texttt{BG-HILBERT-THM-015} \\
Status & \texttt{proved} \\
Lean declaration & \path{ProjectKoios.BandGauge.Hilbert.rankTwoSpectralFieldLossSq_eq_rankTwoSpectralLossSq} \\
Module & \path{ProjectKoios/BandGauge/Hilbert/SpectralLossBridge.lean} \\
\bottomrule
\end{tabularx}

\section{Exact statement}
Under the assumptions of BG-HILBERT-THM-014, if the centered two-state field has
weighted power $P(d)=\kappa$, then its source-normalized residual loss equals the
existing reduced scalar spectral loss:
\[
 L_{\mathrm{field}}^2(s,\lambda)
 =s^2+\kappa(\lambda-1)^2
 =L_{\mathrm{spec}}^2(s,\lambda).
\]

\section{Source and proof}
This theorem \texttt{formalizes} the final spectral identity in Supplementary Section
S1.1 and the exact spectral-loss formula near line 663 of the Paper A v3 manuscript.
Owner: \path{ksdft2effmass}; immutable commit
\begin{center}\small\texttt{99a46cb59d0c6f27ada7d709f194996ef9619d9f}.\end{center}
Substitute $P(d)=\kappa$ into BG-HILBERT-THM-014 and unfold the repository's reduced
scalar spectral-loss definition.

\section{Boundary and trusted base}
This theorem bridges two repository definitions under explicit abstract hypotheses. It
does not establish those hypotheses for scientific source data or certify numerical
results. Lean and pinned Mathlib report \texttt{propext}, \texttt{Classical.choice}, and
\texttt{Quot.sound}. No admission or scientific axiom is used. The proof and prose were
AI-assisted and kernel-checked. No independent expert review is recorded. Status:
\texttt{proved}; working draft.
\end{document}
